\chapter{The equations of the standard model}\label{chap:equations}

The equations are grouped in five blocks: prices (P), production (Q), trade
(T), institutions (I) and the system constraints (S). Their labels are those
that \manref{cge_solve} displays in its diagnostics and that the GAMS program
written by \manref{cge_gams} uses, with a prefix \cmd{E} (\cmd{EP1} is
equation~\eqref{eq:P1}). Each equation holds for every element of the sets
shown with it; a term that refers to a variable that does not exist for an
element (for example imports of a commodity that is not imported) is zero.

Several equations include terms of the extensions: margins by market, VAT,
taxes on factors, rents of constraints, public debt, home consumption and
export demand. Without the extension, the terms are zero (or one, for a
scaling factor) and the equation reduces to its standard form.

\section{Prices}\label{sec:prices}

The import price is the world price, with the tariff and any import premium,
converted at the exchange rate, plus the margins paid per unit of imports:
\begin{equation}
\V{PM}_c = \V{pwm}_c\,(1 + \V{tm}_c\,\V{TMADJ} + \V{RM}_c)\,\V{EXR}
  + \sum_{c'} \V{PQ}_{c'}\,\V{ici}_{c'c}, \qquad c \in CM
\tag{P1}\label{eq:P1}
\end{equation}
The export price received by producers is the world price, less the export
tax and any export premium, less the margins paid per unit of exports:
\begin{equation}
\V{PE}_c = \V{PWE}_c\,(1 - \V{te}_c - \V{RE}_c)\,\V{EXR}
  - \sum_{c'} \V{PQ}_{c'}\,\V{ice}_{c'c}, \qquad c \in CE
\tag{P2}\label{eq:P2}
\end{equation}
with $\V{PWE}_c = \V{pwe}_c$ for $c \notin CED$. The premiums $\V{RM}_c$ and
$\V{RE}_c$ are zero unless a quota or another constraint frees them
(\secref{sec:rents}).

The value of composite supply, net of the sales tax, is the value of domestic
sales at the price domestic buyers pay, of imports, and of the margin
services that composite supply carries:
\begin{equation}
\V{PQ}_c\,(1 - \V{tq}_c\,\V{TQADJ})\,\V{QQ}_c = \V{PDD}_c\,\V{QD}_c
  + \V{PM}_c\,\V{QM}_c + \sum_{c' \in CT} \V{PQ}_{c'}\,\V{icm}_{c'c}\,\V{QQ}_c
\tag{P3}\label{eq:P3}
\end{equation}
where
\begin{equation*}
\V{PDD}_c = \V{PD}_c + \sum_{c'} \V{PQ}_{c'}\,\V{icd}_{c'c}
\end{equation*}
is the price that domestic buyers pay for domestic output, including the
margins on domestic sales ($\V{PDD}_c = \V{PD}_c$ without them). The domestic
term applies only for $c \in CX$, the import term only for $c \in CM$, and the
margin term only when $c$ pays margins.

The value of commodity output is the value of its domestic sales and exports:
\begin{equation}
\V{PX}_c\,\V{QX}_c = \V{PD}_c\,\V{QD}_c + \V{PE}_c\,\V{QE}_c, \qquad c \in CX
\tag{P4}\label{eq:P4}
\end{equation}
where the export term applies only for $c \in CE$. The activity price is the
value of the commodities a unit of activity yields:
\begin{equation}
\V{PA}_a = \sum_c \theta_{ac}\,\V{PX}_c
\tag{P5}\label{eq:P5}
\end{equation}
(for commodities whose output is aggregated by a CES function, $\V{PX}_c$ is
replaced by the activity's own price $\V{PXAC}_{ac}$; \secref{sec:outputces}).
The price of aggregate intermediate input includes the VAT that the activity
cannot recover:
\begin{equation}
\V{PINTA}_a = \sum_c \V{PQ}_c\,\bigl(1 + \V{tv}_c\,\V{TVADJ}\,(1 - \kappa_a)\bigr)\,\V{ica}_{ca}
\tag{P6}\label{eq:P6}
\end{equation}
Revenue net of the activity tax and of any rent on output equals the cost of
value added and intermediate input:
\begin{equation}
\V{PA}_a\,(1 - \V{ta}_a - \V{RA}_a)\,\V{QA}_a = \V{PVA}_a\,\V{QVA}_a + \V{PINTA}_a\,\V{QINTA}_a
\tag{P7}\label{eq:P7}
\end{equation}
The rent $\V{RA}_a$ is zero unless a capacity limit or another constraint
frees it. The price indexes are
\begin{equation}
\V{CPI} = \sum_c \V{cwts}_c\,\V{PH}_c
\tag{P8}\label{eq:P8}
\end{equation}
\begin{equation}
\V{DPI} = \sum_c \V{dwts}_c\,\V{PD}_c
\tag{P9}\label{eq:P9}
\end{equation}
where
\begin{equation*}
\V{PH}_c = \V{PQ}_c\,(1 + \V{tv}_c\,\V{TVADJ})
\end{equation*}
is the price that final users pay, including VAT ($\V{PH}_c = \V{PQ}_c$ without
VAT). $\V{PH}_c$ is an expression, not a variable.

\section{Production}\label{sec:production}

\subsection{The top nest}

An activity combines aggregate value added and aggregate intermediate input.
With a Leontief top nest ($\sigma^a_a = 0$, the default), each is a fixed
quantity per unit of activity:
\begin{equation}
\V{QVA}_a = \V{iva}_a\,\V{QA}_a
\tag{Q1}\label{eq:Q1}
\end{equation}
\begin{equation}
\V{QINTA}_a = \V{inta}_a\,\V{QA}_a
\tag{Q2}\label{eq:Q2}
\end{equation}
With a CES top nest ($\sigma^a_a > 0$, $\sigma^a_a \ne 1$), they substitute
for each other as their prices change:
\begin{equation}
\V{QA}_a = \V{aa}_a\Bigl[\delta^a_a\,\V{QVA}_a^{-\rho^a_a} + (1 - \delta^a_a)\,\V{QINTA}_a^{-\rho^a_a}\Bigr]^{-1/\rho^a_a}
\tag{Q1C}\label{eq:Q1C}
\end{equation}
\begin{equation}
\frac{\V{QVA}_a}{\V{QINTA}_a} = \left[\frac{\V{PINTA}_a}{\V{PVA}_a}\cdot\frac{\delta^a_a}{1 - \delta^a_a}\right]^{1/(1 + \rho^a_a)}
\tag{Q2C}\label{eq:Q2C}
\end{equation}
and with a Cobb--Douglas top nest ($\sigma^a_a = 1$),
\begin{equation}
\V{QA}_a = \V{aa}_a\,\V{QVA}_a^{\delta^a_a}\,\V{QINTA}_a^{1 - \delta^a_a},
\qquad \V{PVA}_a\,\V{QVA}_a\,(1 - \delta^a_a) = \V{PINTA}_a\,\V{QINTA}_a\,\delta^a_a
\tag{Q1C$'$, Q2C$'$}\label{eq:Q1CD}
\end{equation}
In every case \eqref{eq:P7} sets the activity price. An activity without
intermediate inputs keeps the Leontief top nest.

\subsection{Value added}

Factors enter value added in efficiency units,
\begin{equation*}
\widetilde{\V{QF}}_{fa} = \V{af}_f\,\V{EFF}_f\,\V{DW}_f\,\V{QF}_{fa},
\end{equation*}
where $\V{af}_f$ is the exogenous efficiency of the factor, $\V{EFF}_f$ its
efficiency from health services \eqref{eq:H2} and $\V{DW}_f$ the days worked
per worker \eqref{eq:H3}; all three are one in the base. Factor quantities
$\V{QF}_{fa}$, supplies $\V{QFS}_f$ and prices $\V{W}_f$ remain in physical
units, so a more efficient factor earns more per unit.

With CES value added ($\sigma^{va}_a \ne 1$),
\begin{equation}
\V{QVA}_a = \V{ad}_a\Bigl[\sum_f \delta^{va}_{fa}\,\widetilde{\V{QF}}_{fa}^{\,-\rho^{va}_a}\Bigr]^{-1/\rho^{va}_a}
\tag{Q3}\label{eq:Q3}
\end{equation}
and the demand for each factor equates its cost to the value of its marginal
product:
\begin{equation}
\V{W}_f\,\V{WFDIST}_{fa}\,(1 + \V{tf}_{fa}) = \V{PVA}_a\,\V{QVA}_a
  \Bigl[\sum_{f'} \delta^{va}_{f'a}\,\widetilde{\V{QF}}_{f'a}^{\,-\rho^{va}_a}\Bigr]^{-1}
  \delta^{va}_{fa}\,\widetilde{\V{QF}}_{fa}^{\,-\rho^{va}_a - 1}\,\V{af}_f\,\V{EFF}_f\,\V{DW}_f
\tag{Q4}\label{eq:Q4}
\end{equation}
where $\V{tf}_{fa}$ is the tax on the use of the factor (zero without factor
taxes; \secref{sec:factortax}). With Cobb--Douglas value added
($\sigma^{va}_a = 1$),
\begin{equation}
\V{QVA}_a = \V{ad}_a \prod_f \widetilde{\V{QF}}_{fa}^{\,\delta^{va}_{fa}},
\qquad \sum_f \delta^{va}_{fa} = 1
\tag{Q3$'$}\label{eq:Q3CD}
\end{equation}
\begin{equation}
\V{W}_f\,\V{WFDIST}_{fa}\,(1 + \V{tf}_{fa})\,\V{QF}_{fa} = \delta^{va}_{fa}\,\V{PVA}_a\,\V{QVA}_a
\tag{Q4$'$}\label{eq:Q4CD}
\end{equation}
Value added may also combine nests of factors and commodities, each with its
own elasticity; equations \eqref{eq:Q3} and \eqref{eq:Q4} then take the form
of \secref{sec:nests}.

\subsection{Health services and factor efficiency}

When some factor's efficiency depends on health services
(\secref{sec:efficiency}), or a health outcome of chapter~\ref{chap:health}
has a non-zero elasticity, the volume of health services bought by
households, the government and the health insurance fund, at base prices, is
\begin{equation}
\V{HS} = \sum_{c \in HC} \Bigl(\sum_h \V{QH}_{ch} + \V{QG}_c + \V{QSHI}_c\Bigr)
\tag{H1}\label{eq:H1}
\end{equation}
and the efficiency of factor $f$ is
\begin{equation}
\V{EFF}_f = \left(\V{HS} / \V{HS}^0\right)^{\eta_f}, \qquad \eta_f > 0
\tag{H2}\label{eq:H2}
\end{equation}
with $\V{EFF}_f = 1$ otherwise. With a fixed population, $\V{HS}/\V{HS}^0$ is
also the ratio of health services per head to its base value; in dynamic runs
$\V{HS}^0$ grows with the population (chapter~\ref{chap:dynamics}).

\subsection{Intermediate input and output}

Intermediate inputs are Leontief within the aggregate:
\begin{equation}
\V{QINT}_{ca} = \V{ica}_{ca}\,\V{QINTA}_a
\tag{Q5}\label{eq:Q5}
\end{equation}
Output of each commodity is the sum of the yields of the activities that
produce it, applied to their marketed output:
\begin{equation}
\V{QX}_c = \sum_a \theta_{ac}\,\Bigl(\V{QA}_a - \sum_h \V{QHA}_{ah}\Bigr), \qquad c \in CX
\tag{Q6}\label{eq:Q6}
\end{equation}
Output that households consume at home ($\V{QHA}_{ah}$, \secref{sec:home})
does not reach the commodity market. The output of a commodity from different
activities is a perfect substitute and sells at one producer price
$\V{PX}_c$, unless it is aggregated by a CES function
(\secref{sec:outputces}).

\section{Trade}\label{sec:trade}

\subsection{Exports}

For commodities exported from domestic output, $c \in CE$, a CET function
transforms output into exports and domestic sales:
\begin{equation}
\V{QX}_c = \V{at}_c\Bigl[\delta^t_c\,\V{QE}_c^{\rho^t_c} + (1 - \delta^t_c)\,\V{QD}_c^{\rho^t_c}\Bigr]^{1/\rho^t_c}
\tag{T1}\label{eq:T1}
\end{equation}
\begin{equation}
\frac{\V{QE}_c}{\V{QD}_c} = \left[\frac{\V{PE}_c}{\V{PD}_c}\cdot\frac{1 - \delta^t_c}{\delta^t_c}\right]^{1/(\rho^t_c - 1)}
\tag{T2}\label{eq:T2}
\end{equation}
Produced commodities that are not exported, $c \in CNE$, are sold at home:
\begin{equation}
\V{QX}_c = \V{QD}_c
\tag{T3}\label{eq:T3}
\end{equation}
and export-only commodities, $c \in CXE$, are exported in full, in place of
\eqref{eq:T1} and \eqref{eq:T2}:
\begin{equation}
\V{QX}_c = \V{QE}_c
\tag{T3E}\label{eq:T3E}
\end{equation}
Exports that face a finite elasticity of world demand, $c \in CED$, fall as
the world price rises:
\begin{equation}
\V{QE}_c = \overline{\V{qe}}_c \left(\frac{\V{PWE}_c}{\V{pwe}_c}\right)^{-\eta^e_c}
\tag{T10}\label{eq:T10}
\end{equation}
where $\V{pwe}_c$ now shifts the demand curve: at the base world price the
world buys $\overline{\V{qe}}_c$ (\secref{sec:exportdemand}).

Some SAMs record exports that domestic output cannot supply: re-exports of
imported goods, exports from inventories (with a matching fall in stocks), or
exports valued at purchaser prices, including the sales taxes and margins on
the exported goods. For such a commodity the CET is not defined, because its
domestic sales would be zero or negative. Its exports, $c \in CR$, are instead
drawn from composite supply, like any other demand, and sold at the composite
price, with elasticity $\sigma^t_c$ with respect to the world price in local
currency:
\begin{equation}
\V{QER}_c = \overline{\V{qer}}_c\left(\frac{\V{PQ}_c}{\V{EXR}\cdot\V{pwe}_c}\right)^{-\sigma^t_c}
\tag{T9}\label{eq:T9}
\end{equation}
The commodity's domestic output, if any, is sold at home \eqref{eq:T3}.

\subsection{Imports and composite supply}

For imported commodities with domestic output, $c \in CM \cap CX$, an
Armington function combines imports and domestic sales:
\begin{equation}
\V{QQ}_c = \V{aq}_c\Bigl[\delta^q_c\,\V{QM}_c^{-\rho^q_c} + (1 - \delta^q_c)\,\V{QD}_c^{-\rho^q_c}\Bigr]^{-1/\rho^q_c}
\tag{T4}\label{eq:T4}
\end{equation}
\begin{equation}
\frac{\V{QM}_c}{\V{QD}_c} = \left[\frac{\V{PDD}_c}{\V{PM}_c}\cdot\frac{\delta^q_c}{1 - \delta^q_c}\right]^{1/(1 + \rho^q_c)}
\tag{T5}\label{eq:T5}
\end{equation}
For commodities that are not imported, $c \in CNM$ (all of which have
domestic output),
\begin{equation}
\V{QQ}_c = \V{aq}_c\,\V{QD}_c
\tag{T6}\label{eq:T6}
\end{equation}
with $\V{aq}_c = \V{QQ}^0_c / \V{QD}^0_c$. With base prices of one, a sales tax
makes composite supply exceed domestic sales in value units,
$\V{QQ}^0_c = \V{QD}^0_c + \SAM(tax_{sal}, c)$, so the coefficient is not one.
For commodities that are imported only, $c \in CM \setminus CX$,
\begin{equation}
\V{QQ}_c = \V{aq}_c\,\V{QM}_c
\tag{T7}\label{eq:T7}
\end{equation}
with $\V{aq}_c = \V{QQ}^0_c / \V{QM}^0_c$, which exceeds one when the commodity
pays a tariff, sales tax or margins.

\subsection{Margins}

Margins are a fixed quantity of services per unit of composite supply, and,
for margin accounts declared by market, per unit of domestic sales, imports
or exports (\secref{sec:margins-market}). Demand for margin services,
$c \in CT$, is
\begin{equation}
\V{QT}_c = \sum_{c'} \bigl(\V{icm}_{cc'}\,\V{QQ}_{c'} + \V{icd}_{cc'}\,\V{QD}_{c'}
  + \V{ici}_{cc'}\,\V{QM}_{c'} + \V{ice}_{cc'}\,\V{QE}_{c'}\bigr)
\tag{T8}\label{eq:T8}
\end{equation}
and their cost enters the prices through \eqref{eq:P1}, \eqref{eq:P2} and
\eqref{eq:P3}. SAMs often record one margin payment per commodity, without a
split between domestic sales, imports and exports; such margins are
attributed to composite supply, as sales taxes are.

\section{Institutions}\label{sec:institutions}

\subsection{Factor income}

The income of a factor is its earnings in every activity, factor income from
abroad, and any rents paid to it:
\begin{equation}
\V{YF}_f = \sum_a \V{W}_f\,\V{WFDIST}_{fa}\,\V{QF}_{fa} + \V{trnsfr}_{f,row}\,\V{EXR} + R_f
\tag{I1}\label{eq:I1}
\end{equation}
Here and below $R_\varrho$ is the sum of the rents paid to the recipient
$\varrho$ (\secref{sec:rents}); it is zero without constraints. Factor income
is distributed, after contributions to the health insurance fund and taxes on
factor income, in fixed shares:
\begin{equation}
\V{YIF}_{if} = \V{shif}_{if}\,(1 - \V{tshi}_f - \V{tfy}_f)\,\V{YF}_f, \quad i \in I \setminus \{shi\};
\qquad \V{YIF}_{shi,f} = \V{tshi}_f\,\V{YF}_f
\tag{I2}\label{eq:I2}
\end{equation}
Without a fund, $\V{tshi}_f = 0$; without a tax on factor income,
$\V{tfy}_f = 0$.

\subsection{Households and enterprises}

The income of a domestic institution is its factor income, transfers from
other domestic institutions out of their income after direct taxes and
savings, transfers from the government and from abroad, interest on the
domestic public debt it holds, and any rents:
\begin{equation}
\begin{split}
\V{YI}_d = {}& \sum_f \V{YIF}_{df} + \sum_{d' \in D} \V{shii}_{dd'}\,(1 - \V{MPS}_{d'})(1 - \V{TINS}_{d'})\,\V{YI}_{d'}\\
&+ \V{trnsfr}_{d,gov}\,\V{CPI} + \V{trnsfr}_{d,row}\,\V{EXR} + \V{hold}_d\,r^d\,\V{DD} + R_d
\end{split}
\tag{I3}\label{eq:I3}
\end{equation}
A household spends a fixed share of its income after direct taxes and
savings on consumption:
\begin{equation}
\V{EH}_h = \V{hc}_h\,(1 - \V{MPS}_h)(1 - \V{TINS}_h)\,\V{YI}_h
\tag{I4}\label{eq:I4}
\end{equation}
Enterprises do not consume: for $e \in E$, $\sum_d \V{shii}_{de} + \V{shir}_e = 1$,
so all their income after direct taxes and savings is transferred. For
households, $\V{hc}_h = 1 - \sum_d \V{shii}_{dh} - \V{shir}_h$. A transfer
from an institution to itself (a diagonal SAM cell) is allowed; it appears on
both sides of its accounts.

Household demand is a linear expenditure system: households buy subsistence
quantities and spend the rest of their consumption budget, their
supernumerary spending
\begin{equation*}
\V{SN}_h = \V{EH}_h - \sum_c \V{PH}_c\,\gamma_{ch} - \sum_a \V{PA}_a\,\gamma^h_{ah},
\end{equation*}
in fixed marginal shares:
\begin{equation}
\V{PH}_c\,\V{QH}_{ch} = \V{PH}_c\,\gamma_{ch} + \beta_{ch}\,\V{SN}_h
\tag{I5}\label{eq:I5}
\end{equation}
\begin{equation}
\V{PA}_a\,\V{QHA}_{ah} = \V{PA}_a\,\gamma^h_{ah} + \beta^h_{ah}\,\V{SN}_h
\tag{I5H}\label{eq:I5H}
\end{equation}
With $\gamma_{ch} = 0$ this is Cobb--Douglas demand, the default: fixed budget
shares and income elasticities of one. Home consumption \eqref{eq:I5H} is
valued at the activity price, without taxes or margins; the marginal shares of
marketed and home goods together sum to one.

\subsection{Investment and government demand}

\begin{equation}
\V{QINV}_c = \overline{\V{qinv}}_c\,\V{IADJ}
\tag{I6}\label{eq:I6}
\end{equation}
\begin{equation}
\V{QG}_c = \overline{\V{qg}}_c\,\V{GADJ}
\tag{I7}\label{eq:I7}
\end{equation}

\subsection{Government}

Government revenue is the direct taxes, the taxes on sales, imports, exports,
activities and factors, net VAT, the government's factor income, transfers
from abroad and within government, and any rents paid to it:
\begin{equation}
\begin{split}
\V{YG} = {}& \sum_d \V{TINS}_d\,\V{YI}_d + \sum_c \V{tq}_c\,\V{TQADJ}\,\V{PQ}_c\,\V{QQ}_c
  + \sum_{c \in CM} \V{tm}_c\,\V{TMADJ}\,\V{pwm}_c\,\V{QM}_c\,\V{EXR}\\
&+ \sum_{c \in CE} \V{te}_c\,\V{PWE}_c\,\V{QE}_c\,\V{EXR} + \sum_a \V{ta}_a\,\V{PA}_a\,\V{QA}_a + \V{VAT}\\
&+ \sum_{f,a} \V{tf}_{fa}\,\V{W}_f\,\V{WFDIST}_{fa}\,\V{QF}_{fa} + \sum_f \V{tfy}_f\,\V{YF}_f
  + \sum_f \V{YIF}_{gov,f}\\
&+ \V{trnsfr}_{gov,row}\,\V{EXR} + \V{trnsfr}_{gov,gov}\,\V{CPI} + R_{gov}
\end{split}
\tag{I8}\label{eq:I8}
\end{equation}
with VAT revenue, net of the credits that activities recover,
\begin{equation*}
\V{VAT} = \V{TVADJ} \sum_c \V{tv}_c\,\V{PQ}_c \Bigl(\sum_h \V{QH}_{ch} + \V{QG}_c + \V{QINV}_c + \V{QSHI}_c
  + \sum_a (1 - \kappa_a)\,\V{QINT}_{ca}\Bigr).
\end{equation*}
Government spending is its consumption, its transfers, and the interest on
public debt:
\begin{equation}
\begin{split}
\V{EG} = {}& \sum_c \V{PH}_c\,\V{QG}_c + \sum_d \V{trnsfr}_{d,gov}\,\V{CPI} + \V{trnsfr}_{gov,gov}\,\V{CPI}
  + \V{trnsfr}_{shi,gov}\,\V{CPI}\\
&+ \V{trnsfr}_{row,gov}\,\V{EXR} + r^d\,\V{DD} + r^f\,\V{DF}\,\V{EXR}
\end{split}
\tag{I9}\label{eq:I9}
\end{equation}
Direct tax and savings rates are their base values times scaling factors that
the closure fixes or frees:
\begin{equation}
\V{TINS}_d = \overline{\V{tins}}_d\,\V{TINSADJ}
\tag{I10}\label{eq:I10}
\end{equation}
\begin{equation}
\V{MPS}_d = \overline{\V{mps}}_d\,\V{MPSADJ}
\tag{I11}\label{eq:I11}
\end{equation}

\subsection{The health insurance fund}

When a social health insurance fund is mapped (\secref{sec:fund}), its
revenue is the contributions on factor income and the transfers from the
government and from abroad:
\begin{equation}
\V{YSHI} = \sum_f \V{YIF}_{shi,f} + \V{trnsfr}_{shi,gov}\,\V{CPI} + \V{trnsfr}_{shi,row}\,\V{EXR}
\tag{I12}\label{eq:I12}
\end{equation}
It buys a fixed basket of health services, scaled so that it spends its
revenue less a fixed share saved:
\begin{equation}
\V{QSHI}_c = \overline{\V{qshi}}_c\,\V{SHIADJ}
\tag{I13}\label{eq:I13}
\end{equation}
\begin{equation}
\sum_c \V{PH}_c\,\V{QSHI}_c = (1 - \V{sshi})\,\V{YSHI}
\tag{I14}\label{eq:I14}
\end{equation}

\section{System constraints}\label{sec:system}

Factor markets clear:
\begin{equation}
\sum_a \V{QF}_{fa} = \V{QFS}_f
\tag{S1}\label{eq:S1}
\end{equation}
Commodity markets clear: composite supply equals intermediate demand,
household, government and investment demand, margin services, exports from
composite supply and the purchases of the health insurance fund:
\begin{equation}
\V{QQ}_c = \sum_a \V{QINT}_{ca} + \sum_h \V{QH}_{ch} + \V{QG}_c + \V{QINV}_c + \V{QT}_c + \V{QER}_c + \V{QSHI}_c
\tag{S2}\label{eq:S2}
\end{equation}
where $\V{QT}_c$ applies only for $c \in CT$ and $\V{QER}_c$ only for
$c \in CR$. The balance of payments, in foreign currency, equates payments
abroad (imports, factor income paid abroad, transfers abroad from domestic
institutions and the government, interest on foreign debt, and rents paid
abroad) to receipts (exports, transfers and factor income from abroad, and
foreign savings):
\begin{equation}
\begin{split}
&\sum_{c \in CM} \V{pwm}_c\,\V{QM}_c + \sum_f \frac{\V{YIF}_{row,f}}{\V{EXR}}
  + \sum_d \frac{\V{shir}_d\,(1 - \V{MPS}_d)(1 - \V{TINS}_d)\,\V{YI}_d}{\V{EXR}}\\
&\qquad + \V{trnsfr}_{row,gov} + r^f\,\V{DF} + \frac{R_{row}}{\V{EXR}}\\
&\quad = \sum_{c \in CE} \V{PWE}_c\,\V{QE}_c + \sum_{c \in CR} \frac{\V{PQ}_c\,\V{QER}_c}{\V{EXR}}
  + \sum_d \V{trnsfr}_{d,row} + \sum_f \V{trnsfr}_{f,row}\\
&\qquad + \V{trnsfr}_{gov,row} + \V{trnsfr}_{shi,row} + \V{FSAV}
\end{split}
\tag{S3}\label{eq:S3}
\end{equation}
The government's budget defines its savings:
\begin{equation}
\V{YG} = \V{EG} + \V{GSAV}
\tag{S4}\label{eq:S4}
\end{equation}
and savings equal investment:
\begin{equation}
\sum_d \V{MPS}_d\,(1 - \V{TINS}_d)\,\V{YI}_d + \V{GSAV} + \V{sshi}\,\V{YSHI} + \V{EXR}\cdot\V{FSAV}
  = \sum_c \V{PH}_c\,\V{QINV}_c + \V{WALRAS}
\tag{S5}\label{eq:S5}
\end{equation}
The price level is set by the numeraire:
\begin{equation}
\V{CPI} = \overline{\V{CPI}} \quad\text{or}\quad \V{DPI} = \overline{\V{DPI}}
\tag{S6}\label{eq:S6}
\end{equation}
which is imposed by fixing the numeraire variable in the closure
(chapter~\ref{chap:closures}); it is not counted as a separate equation.

By Walras's law one equation is redundant: if every other market clears and
every institution meets its budget, savings equal investment. The variable
$\V{WALRAS}$ makes the system square, and must be zero, to the solver's
tolerance, at every solution. \manref{cge_results} and \manref{cge_check}
report it.
